%! Skills Focus: Selecting, Implementing, and Communicating Inference Procedures — Unit 7, Lesson 7.10 — Group Activity (Answer Key)
\documentclass[10pt]{article}
\usepackage{apstats-article}
\usepackage{apstats-key}   % pulls in -boxes; do NOT also load -boxes
\newcommand{\TallMath}[1]{$\displaystyle #1\rule[-1.4em]{0pt}{3.2em}$}

\begin{document}

\pageheader{Unit 7, Lesson 7.10}{Group Activity --- Answer Key}
\namepartnerperiod

\vspace{0.10in}

\textbf{Procedure Menu} (use for Tier R): one-sample $z$-interval for $p$ \;|\;
one-sample $z$-test for $p$ \;|\; two-sample $z$-test for $p_1-p_2$ \;|\;
one-sample $t$-test for $\mu$ \;|\; two-sample $t$-test for $\mu_1-\mu_2$ \;|\;
matched-pairs $t$-test for $\mu_d$

\vspace{0.08in}

\begin{scenariobox}[Scenario 1 --- Sleep and Grade Level]{sky}
A high school counselor randomly selects 25 freshmen and 25 seniors and records the
number of hours each student slept the previous night. Summary statistics:
Freshmen: $\bar x_1 = 7.1$ hr, $s_1 = 0.9$ hr. \;
Seniors: $\bar x_2 = 6.4$ hr, $s_2 = 1.1$ hr.
The dotplots for both groups show roughly symmetric distributions with no outliers.
She wants to test whether freshmen sleep more, on average, than seniors.

\textbf{Tier R:} Select the procedure from the menu above and list the conditions.
\begin{itemize}
  \item Procedure: \ans{two-sample $t$-test for $\mu_1-\mu_2$}
  \item Conditions: Random --- \ans{random selection} \; 10\% --- \ans{$25 \le 10\%$ of all students} \; Normal --- \ans{no outliers; symmetric}
\end{itemize}

\textbf{Tier A:} Complete all four steps.
\begin{enumerate}
  \item \textbf{State:} Procedure name: \ans{two-sample $t$-test for $\mu_1 - \mu_2$}\\
        $H_0$: \ans{$\mu_1 - \mu_2 = 0$} \quad $H_a$: \ans{$\mu_1 - \mu_2 > 0$}
  \item \textbf{Plan:} Verify conditions (with numbers):
        \begin{itemize}
          \item Random: \ansline{Both groups were randomly selected --- met.}
          \item 10\%: \ansline{$25 \le 10\%$ of all freshmen/seniors at the school (assuming enrollment $> 250$) --- met.}
          \item Normal: \ansline{$n_1 = n_2 = 25 < 30$; dotplots show symmetric distributions with no outliers --- Normal condition met.}
        \end{itemize}
  \item \textbf{Do:} $t = \dfrac{(\bar x_1 - \bar x_2)}{SE}$ \quad
        $SE = \sqrt{\dfrac{s_1^2}{n_1} + \dfrac{s_2^2}{n_2}} = \sqrt{\dfrac{0.9^2}{25} + \dfrac{1.1^2}{25}} \approx$ \ans{0.284}

        $t \approx$ \ans{2.46} \quad $df \approx$ \ans{46} (technology) \quad $p$-value $\approx$ \ans{0.009}
  \item \textbf{Conclude:} \ansline{Because $p \approx 0.009 < \alpha = 0.05$, reject $H_0$. There is sufficient evidence that freshmen sleep more, on average, than seniors.}
\end{enumerate}
\end{scenariobox}

\begin{scenariobox}[Scenario 2 --- Vaccine Side Effects]{sky}
A pharmaceutical company randomly samples 200 participants from a clinical trial.
Of 100 who received the vaccine, 18 reported a side effect. Of 100 who received a
placebo, 9 reported a side effect. The company wants to estimate how much higher the
proportion with side effects is for the vaccine group compared to the placebo group.

\textbf{Tier R:} Select the procedure and list the conditions.
\begin{itemize}
  \item Procedure: \ans{two-sample $z$-interval for $p_1 - p_2$}
  \item Large Counts (vaccine): $n_1 \hat p_1 =$ \ans{18} and $n_1(1-\hat p_1) =$ \ans{82}
  \item Large Counts (placebo): $n_2 \hat p_2 =$ \ans{9} and $n_2(1-\hat p_2) =$ \ans{91}
\end{itemize}

\textbf{Tier A:} Construct a 95\% confidence interval.
\begin{enumerate}
  \item \textbf{State:} Procedure: \ans{two-sample $z$-interval for $p_1-p_2$} \quad
        Parameter: $p_1 - p_2 =$ \ans{true difference in proportions reporting a side effect (vaccine $-$ placebo)}
  \item \textbf{Plan:} Verify conditions (with numbers):
        \begin{itemize}
          \item Random: \ansline{Randomly sampled participants --- met.}
          \item 10\%: \ansline{$100 \le 10\%$ of all trial participants (assumed large) --- met.}
          \item Large Counts: \ansline{All four counts $\ge 10$: 18, 82, 9, 91 --- met.}
        \end{itemize}
  \item \textbf{Do:} $\hat p_1 =$ \ans{0.18} \; $\hat p_2 =$ \ans{0.09} \;
        $SE = \sqrt{\dfrac{0.18(0.82)}{100} + \dfrac{0.09(0.91)}{100}} \approx$ \ans{0.0459}

        Interval: $(0.18-0.09) \pm 1.96 \times$ \ans{0.0459} $= ($ \ans{$-0.001$}, \ans{$0.181$} $)$
  \item \textbf{Conclude:} \ansline{We are 95\% confident that the proportion reporting side effects is between $-0.001$ and $0.181$ higher for the vaccine group than the placebo group. Because 0 is just inside the interval, this is not conclusive evidence of a higher rate.}
\end{enumerate}
\end{scenariobox}

\begin{scenariobox}[Scenario 3 --- Tier E: Critique a Student Response]{sky}
\textbf{Tier E only.} \; A student was given this prompt: ``A researcher measures reaction time
(milliseconds) for the same 15 subjects before and after caffeine. He finds
$\bar d = -12$ ms, $s_d = 8$ ms, and wants to test whether caffeine reduces mean reaction time.''

The student wrote:
\begin{quote}
\textit{``I will use a two-sample $t$-test for $\mu_1 - \mu_2$. $H_0: \bar x_1 = \bar x_2$;
$H_a: \bar x_1 > \bar x_2$. Random: the 15 subjects were selected. Normal: $n = 15 < 30$ so
the condition fails. Therefore we cannot proceed.''
}
\end{quote}

Identify \textbf{two} errors in the student's response and write a corrected version of the
affected steps.

\begin{enumerate}
  \item Error 1: \ansline{Wrong procedure: same subjects measured twice = matched pairs, not two independent samples. A two-sample $t$-test is inappropriate.}
        Correction: \ansline{Use a matched-pairs $t$-test for $\mu_d$, where $\mu_d$ = mean difference in reaction time (before $-$ after caffeine) for all subjects.}
  \item Error 2: \ansline{Two sub-errors: (a) hypotheses use sample statistics ($\bar x_1, \bar x_2$) instead of parameters; (b) the Normal condition conclusion is wrong --- $n < 30$ does not automatically mean the condition fails.}
        Correction: \ansline{Hypotheses: $H_0: \mu_d = 0$; $H_a: \mu_d > 0$ (if before $-$ after, a positive difference means caffeine reduced reaction time). Normal: $n = 15 < 30$; check the dotplot of differences for strong skewness or outliers. If none present, proceed.}
\end{enumerate}

\begin{teachernote}
Tier E students should be able to explain why matched pairs is correct (dependence between
measurements on the same subject) and articulate why $n < 30$ alone does not fail the
Normal condition for $t$-procedures.
\end{teachernote}
\end{scenariobox}

\end{document}
